\subsection{关于子空间的极限}

设 X、Y 是两个拓扑空间，A 是 X 的子集，且 \(a \in X\) 是 A 的闭包中的一点（但不一定属于 A）。设 F 是 a 在 X 中的邻域滤子在 A 上的迹。若 f 是 A 到 Y 中的映射，则我们不把“\(y \in Y\) 是 f 关于 F 的极限”写作 \(y = \lim_{\mathfrak{F}} f\)，而记作

\[
y = \lim _ {x \to a, x \in \mathbf {A}} f (x)
\]

并称 \(y\) 为 \(f\) 在 \(a\) 处关于子空间 \(A\) 的极限，或称 \(f(x)\) 趋于 \(y\)（当 \(x\) 趋于 a 且保持在 \(A\) 中时）。此时有 \(y \in \overline{f(A)}\)。

若 \(\mathbf{A} = \complement\{a\}\)，其中 \(a\) 不是 \(\mathbf{X}\) 的孤立点，则我们用 \(y = \lim_{x\to a,x\neq a}f(x)\) 代替 \(y = \lim_{x\to a,x\in \mathbf{A}}f(x)\)。

对聚点可作类似的定义。

若 \(f\) 是 \(\mathbf{A}\) 上的限制（映射 \(g: \mathbf{X} \to \mathbf{Y}\) 的），则称 \(g\) 有极限（相应地，聚点）\(y\)（关于 \(\mathbf{A}\)，在点 \(a \in \overline{\mathbf{A}}\) 处），如果 \(y\) 是 \(f\) 在 \(a\) 处关于 \(\mathbf{A}\) 的极限（相应地，聚点）。

设 B 是 A 的子集，且 \(a \in X\) 是 B 的闭包中的一点；若 y 是映射 \(f: A \to Y\) 在 a 处关于 A 的极限，则 y 也是 f 在 a 处关于 B 的极限；反之未必成立。但若 V 是点 \(a \in \overline{A}\) 在 X 中的一个邻域，且 f 在 a 处关于 \(V \cap A\) 有极限 y，则 y 仍是 f 在 a 处关于 A 的极限。

设 \(a\) 是 \(\mathbf{X}\) 的非孤立点，于是 \(a\) 属于 \(\complement\{a\}\) 的闭包。则映射 \(f: \mathbf{X} \to \mathbf{Y}\) 在 \(a\) 连续，当且仅当 \(f(a) = \lim_{x \to a, x \neq a} f(x)\)；这由定义立得。

